\subsection*{Question order in visual reading}
In the previous experiment the queried glyph was marked. Here the reader received an unmarked clause and a pictured question asking for the English meaning of its main verb or direct object. On 539 held clauses, a reader trained from the answer images located and glossed both terms on 225 when asked in previously unseen Chinese wording. Providing locations during training also gave 225. Exchanging the two training answers within pages gave zero, while a local training-count system supplied with the role identity completed 230. Giving the writer both correct source locations raised its own score to 363.

We next asked for one grammatical role while explicitly excluding the other. The two requests contain the same glyphs and every three-cell window occurs equally often. The recurrent reader selected the correct source glyph and wrote its meaning on every request across both word orders and both languages for 230/539 clauses (Fig.~\ref{fig:asked-gloss}). The matched pooled architecture and the recurrent exchanged-answer control each scored zero. In the excerpt about holding the great image, the released model produced \textit{hold} for the verb and \textit{image} for the object. The reader's success depended on question wording. Four later pairs avoiding the terms for verb and object gave 0, 33, 0 and 2 complete pages respectively. Only 13 held clauses placed the object before the verb. The ordered reader completed one of them under the positive-first Chinese contrast. The held sources contained no cross-role contrasts with an identical three-cell neighbourhood under the prescribed pairing rule, which leaves long-range role assignment unresolved.

An optional decoder replaced 505/3,119 freely produced source-position streams that fell outside the physical training glossary or failed to stop. It recovered 44 accepted readings at the true positions and lost none on this sample. Complete eight-request scores rose to 247/539. On the separately sourced \textit{Dao De Jing}, the increase was from four to five of 17. One erroneous letter sequence for the male term became \textit{male}; the sequence for the female term became \textit{turn}. The lexical constraint can enforce an attested spelling even where the chosen meaning remains wrong.
